\documentclass[12pt,a4paper]{article}
\usepackage[utf8]{inputenc}
\usepackage[margin=2.5cm]{geometry}
\usepackage{amsmath,amssymb,amsfonts}
\usepackage{xcolor}
\usepackage{url}
\usepackage{graphicx}
\usepackage{booktabs}
\usepackage[hidelinks]{hyperref}
\usepackage{setspace}
\onehalfspacing

% Color definitions for changes
\definecolor{revcolor}{RGB}{0,0,180}
\newcommand{\rev}[1]{\textcolor{revcolor}{#1}}
\newcommand{\RC}[1]{\noindent\textit{#1}\vspace{0.3em}\\}
\newcommand{\AR}{\noindent\textbf{Response:}~}

\begin{document}

\begin{center}
{\Large \textbf{Response to Reviewers}}\\[0.8em]
{\normalsize Manuscript ID: CAIS-D-25-01702}\\[0.3em]
{\normalsize \textit{Byzantine-Resilient Stochastic Games for Federated Multi-Cloud Intrusion Detection}}\\[0.3em]
{\normalsize Roger Nick Anaedevha, Alexander G. Trofimov, Yuri V. Borodachev}\\[0.3em]
{\normalsize National Research Nuclear University MEPhI, Moscow, 115409, Russia}
\end{center}

\vspace{1em}

\noindent Dear Editor-in-Chief Professor Yaochu Jin,

\vspace{0.5em}

\noindent We sincerely thank you and the reviewers for the constructive feedback on our manuscript. We have carefully addressed every comment and carried out a thorough revision that strengthens both the theoretical and empirical contributions of the paper. Below we provide a point-by-point response to each reviewer's comments. Throughout the revised manuscript, all substantive changes appear in \rev{blue text} for easy identification. We have also condensed the writing, removed unnecessary verbosity, and restructured several sections to improve clarity and readability, as suggested.

\vspace{1.5em}
\hrule
\vspace{1em}

%% ============================
\section*{Response to Reviewer \#1}
%% ============================

\RC{Comment 1.1: The game-theoretic part is introduced as one that provides formal guarantees of equilibrium, while the manuscript does not provide a good enough explanation of what the structure of the game is, what the utility functions of the players are, or what type of equilibrium concept has been considered.}

\AR We appreciate this important observation. In the revised manuscript we have substantially expanded Section~4.1 to make the game structure explicit. We now clearly state that the game belongs to the class of multi-player, non-zero-sum stochastic games with incomplete information, played between a coalition of $K$ defenders and an adversary coalition of $M$ attackers. The strategy spaces for each player type are formally defined: for each defender~$D_k$ in domain~$d$, the action set $\mathcal{A}_k^{(d)}$ consists of detection threshold selections, model parameter updates, and resource allocation decisions; for the adversary, the action set $\mathcal{A}_A$ comprises domain-specific attack vectors together with cross-domain coordination choices. The utility functions are now presented in Definition~4.3, where each defender's payoff $u_k^{(d)}(s,a)$ is a weighted sum of detection rewards, false-positive costs, and resource expenditures, with class-specific inverse-frequency weights $w_c^{(d)}$ that correct for extreme imbalance. The adversary objective $J_A$ (Equation~3) is a time-integrated expected return incorporating infiltration success, cross-domain synergies, detection penalties, and resource constraints. We adopt a mixed-strategy Nash equilibrium concept, where no player can unilaterally improve their expected payoff by changing strategy. Theorem~4.4 proves existence of this equilibrium via Kakutani's fixed-point theorem and provides an imbalance-robustness bound quantifying the deviation from uniform strategies. The revised Section~4.1.1 also includes an intuitive walk-through of the game structure before the formal definitions, so that readers can follow the mathematical formalization more easily.

These changes appear in Sections~4.1.1 through 4.1.4 of the revised manuscript (pp.~13--16).

\vspace{0.8em}

\RC{Comment 1.2: Without explicit mathematical formulation and proof sketches, the claimed theoretical guarantees are not easy to judge.}

\AR We have addressed this concern in two ways. First, the main manuscript now includes complete proof sketches for all four main theorems (Theorems~4.2, 4.4, 4.5, and 4.7), each outlining the key technical steps and the mathematical tools employed. Second, the revised Supplementary Material (Appendix~C) provides full, self-contained proofs. In particular, the proof of Theorem~4.2 (existence and uniqueness) relies on Picard iteration in the space of c\`adl\`ag processes with Gr\"onwall's inequality; the proof of Theorem~4.4 (Nash equilibrium existence) applies Kakutani's fixed-point theorem to the best-response correspondence with imbalance-weighted payoffs; the proof of Theorem~4.5 (privacy and robustness) decomposes the guarantee into a sensitivity analysis via gradient clipping and a Gaussian mechanism analysis; and the proof of Theorem~4.7 (almost-sure convergence) constructs a supermartingale from the heterogeneous Lyapunov function and applies the Robbins--Siegmund theorem. Each proof now clearly states the assumptions required, the mathematical tools used, and the intermediate lemmas invoked.

These changes appear in Section~4 of the main manuscript and Appendix~C of the Supplementary Material.

\vspace{0.8em}

\RC{Comment 1.3: The strategic model should be clearly defined, how it integrates into the federated optimization process should be explained, and whether the equilibrium conditions are obtained analytically or empirically validated.}

\AR Thank you for raising this point. In the revised manuscript, we have added a new subsection (Section~4.1.5, ``Integration of Game-Theoretic and Federated Learning Components'') that explicitly describes how the strategic model connects with the federated optimization loop. At each federated round $t$, the defender strategies are updated through local model training and secure aggregation (Algorithm~2, lines 14--22), while the adversary strategy is updated based on observed cross-domain detection rates (line~28). The Nash gap $\Delta_{\text{Nash}}^{(t)}$ (line~29) measures distance from equilibrium and serves as the convergence criterion. Regarding the nature of the equilibrium, we clarify that our approach is a hybrid: the existence of Nash equilibrium is established analytically through Theorem~4.4, while convergence to this equilibrium under the federated learning dynamics is proved analytically in Theorem~4.7. We then empirically validate the convergence rate by tracking the Nash gap across federated rounds, as shown in the new Figure~6 in Section~6.5, which plots the empirical Nash gap alongside the theoretical $O(T^{-2/3})$ bound.

\vspace{0.8em}

\RC{Comment 1.4: Your paper does not fully explain which aggregation rule is used, how class imbalance severity is dealt with across the domains, and how feature heterogeneity is reconciled among the IoT, containerized systems, and SOC data streams.}

\AR We have expanded Section~4.2 to address each of these points in detail. Regarding the aggregation rule, we now provide a step-by-step explanation of our hierarchical aggregation procedure (Algorithm~1): within each domain, gradients are first projected onto a common subspace via domain-specific projection matrices~$P_d$, then filtered through cosine-similarity-based Byzantine detection, aggregated using a coordinate-wise trimmed mean with domain-specific trim ratios~$\alpha_d$, and finally combined across domains through imbalance-weighted averaging. The trimmed mean removes the top and bottom $\alpha_d/2$ fraction of values for each coordinate, where $\alpha_d$ is set to twice the expected Byzantine fraction per domain. For class imbalance, we employ inverse square-root frequency weighting $w_c^{(d)} = \sqrt{n_{\text{total}}^{(d)}/(C_d \cdot n_c^{(d)})}$ applied at three levels: local loss computation, gradient clipping norms, and cross-domain aggregation weights. This triple-level weighting ensures that minority classes contribute meaningfully at each stage of the federated pipeline. For feature heterogeneity, we now include a detailed description of the projection mechanism: each domain maintains domain-specific input layers that map raw features to a shared latent representation of dimension 128, and only the shared layers are aggregated across domains. Domain-specific layers remain local to preserve the unique characteristics of each security environment.

These changes appear in Sections~4.2 and 3.3 of the revised manuscript.

\vspace{0.8em}

\RC{Comment 1.5: More papers need to be reviewed related to security, IoT and AI [five specific papers listed].}

\AR We thank the reviewer for these suggestions. We have carefully reviewed each paper and incorporated them into our revised manuscript where they materially support or extend our discussion.

Hezam et al.~[26] is cited in the revised Section~2.1 in the context of multi-criteria decision-making frameworks for evaluating security and privacy in emerging technology environments, connecting their spherical fuzzy evaluation methodology to our multi-domain security assessment approach. Abdel-Basset et al.~[27] is cited in Section~3.2 where we discuss threat models for autonomous and cyber-physical systems, drawing parallels between their security-by-design risk management framework and our game-theoretic defense modeling. Ding et al.~[28] is cited in the revised Section~2.2 when discussing intelligent forensics techniques for IoT communications and their relevance to cross-domain threat detection in heterogeneous environments. Ding et al.~[29] is cited in Section~7.3 (Future Directions) in the context of emerging roles of large language models in cyber resilience and their potential integration with federated security frameworks. Sallam et al.~[30] is cited in Section~7.2 in the discussion of AI-enhanced service frameworks for smart city environments and their security implications for multi-cloud deployments.

These references now appear as [26]--[30] in the revised bibliography.

\vspace{0.8em}

\RC{Comment 1.6: More algorithm descriptions and ablation studies on individual robustness mechanisms and comparison with the existing robust federated aggregation methods will enhance the credibility.}

\AR We have substantially expanded the ablation study in Section~6.2 (Table~3 in the revised manuscript). Beyond the component-level ablation already present, we now include a dedicated ablation on individual robustness mechanisms in a new Table~8 (Section~6.3), which isolates the contribution of each Byzantine defense component: gradient similarity screening alone, trimmed mean aggregation alone, reputation tracking alone, and their combinations. We also added a direct comparison of our aggregation rule against five established robust aggregation methods (Krum, Multi-Krum, coordinate-wise median, trimmed mean, and BALANCE) under identical experimental conditions, presented in the new Table~9. These results demonstrate that our integrated approach outperforms each individual mechanism and each standalone robust aggregation method across all three domains.

\vspace{0.8em}

\RC{Comment 1.7: The experimental check is enormous in scale... The composition of the dataset, labeling mechanism, and cross-domain partitioning need to be described.}

\AR We have expanded Section~5.1 to include a comprehensive description of the dataset composition, labeling procedures, and cross-domain partitioning strategy. For the Edge-IIoT component, we now describe the labeling mechanism: network flows are labeled using a combination of ground-truth attack injection logs (for simulated attacks) and Snort/Suricata IDS rule matching (for protocol-level classification). For the Container component, labels correspond to specific CVE exploitation timestamps recorded during controlled Kubernetes penetration testing. For the SOC component, labels are derived from Microsoft Security Analyst annotations reflecting real-world incident triage decisions. Regarding cross-domain partitioning, we describe our Dirichlet-based allocation procedure (with concentration parameter $\alpha = 0.3$) and explain how we ensure that the temporal ordering of events is preserved within each client's partition while maintaining realistic heterogeneity across clients.

\vspace{0.8em}

\RC{Comment 1.8: The robustness guarantees for Byzantine corruption are compelling... However, there are assumptions in terms of threat models regarding the knowledge, rate, or collusions of the adversaries.}

\AR This is an important clarification. We have added a new Section~3.2.1 (``Threat Model Assumptions and Boundaries'') that explicitly lists all assumptions underlying our Byzantine resilience guarantees. Specifically, we assume that Byzantine clients constitute at most $f_d < \alpha_d/2$ of the clients in any single domain, that Byzantine clients can send arbitrary gradient updates but cannot observe other clients' private gradients before submitting their own, and that the honest clients' gradients satisfy a bounded variance condition. We also state that our threat model does not assume knowledge of whether adversaries collude, but our experimental evaluation in Section~6.3 explicitly tests coordinated attacks where all Byzantine clients simultaneously execute the same attack strategy. The revised Table~4 now includes results for both independent and colluding adversaries, showing that the framework retains 92.8\% performance under coordinated attacks at 20\% corruption, compared to 94.0\% average across independent strategies.

\vspace{0.8em}

\RC{Comment 1.9: The claim of scalability to millions of events per second on major cloud providers is certainly ambitious and impactful, but the description of the deployment does not include architectural detail.}

\AR We have expanded Section~6.7 and added a new deployment architecture description in Section~4.5 that provides specific architectural details. The revised text now describes the three-tier deployment: an ingestion layer using Apache Kafka for event streaming with domain-specific topic partitioning, a processing layer employing PyTorch-based model inference on NVIDIA A100 GPUs with batch-optimized data loaders, and an aggregation layer implementing our secure federated protocol over gRPC with TLS 1.3 encryption. We also report per-component latency breakdowns in the revised Table~7, showing that feature extraction accounts for 42\% of inference time for Edge-IIoT, model forward pass accounts for 31\%, and post-processing accounts for 27\%. The aggregate throughput of 3.5 million events per second is achieved through horizontal scaling across four GPU-equipped nodes, each handling domain-specific inference in parallel.

\vspace{0.8em}

\RC{Comment 1.10: Using your language and drawing conclusions more directly from the results is more appropriate for academic tone.}

\AR We have revised the writing throughout the manuscript to adopt a more measured academic tone. We replaced promotional language such as ``remarkable,'' ``exceptional,'' and ``ambitious'' with factual, data-driven statements. For example, ``demonstrates exceptional Byzantine resilience'' has been changed to ``retains 94.0\% average performance under 20\% Byzantine corruption.'' We have also restructured the Discussion section to draw conclusions more directly from the experimental observations rather than restating claims, and we have reduced redundancy between the Results and Discussion sections.

\vspace{0.8em}

\RC{Comment 1.11: Categorizing theoretical, algorithmic, and system engineering advancements under appropriately defined sections would further improve clarity.}

\AR We have restructured Section~4 (Methodology) into three clearly delineated subsections: Section~4.1 covers the theoretical game-theoretic framework (Definitions~4.1--4.3, Theorems~4.2 and 4.4), Section~4.2 addresses the algorithmic contributions including the Byzantine-resilient protocol (Algorithm~1, Theorem~4.5) and convergence analysis (Theorem~4.7), and Section~4.4--4.5 describe the system engineering aspects including domain-adaptive adversarial training and overall system architecture. The revised Section~1.3 (Research Contributions) now explicitly maps each contribution to its corresponding section, making it easier for readers to locate the relevant technical details.

\vspace{1.5em}
\hrule
\vspace{1em}

%% ============================
\section*{Response to Reviewer \#2}
%% ============================

\RC{Comment 2.1: The payoff functions, strategy spaces, and sensitivity analysis of the game theoretic model may be provided.}

\AR We have addressed this in the revised manuscript as follows. The payoff functions are now explicitly defined in Definition~4.3 (Equation~10), where each defender's utility is a class-weighted sum of detection rewards, false-positive costs, and resource expenditures. The strategy spaces are formally specified in the new Definition~4.1: each defender selects from a continuous action space comprising detection threshold parameters and model update directions, while the adversary selects from a mixed-strategy space over domain-specific attack vectors.

Regarding sensitivity analysis, we have added a new Section~6.8 (``Game-Theoretic Sensitivity Analysis'') that examines how the Nash equilibrium outcomes respond to perturbations in key model parameters. We conducted a one-at-a-time parametric analysis varying each parameter by $\pm 25\%$ around its baseline value, covering the defender detection reward coefficient~$\alpha^{(d)}$, false-positive cost coefficient~$\beta^{(d)}$, adversary cross-domain synergy parameter~$\gamma$, discount factor~$\gamma_{\text{disc}}$, and imbalance ratios~$\rho_d$. The results, presented in the new Figure~8 and Table~10, show that the equilibrium detection accuracy is most sensitive to the false-positive cost coefficient (a 25\% increase in $\beta^{(d)}$ reduces detection accuracy by 2.1\% on average) and least sensitive to the discount factor. The Nash equilibrium remains stable across all tested perturbation ranges, with maximum deviation of 3.4\% from baseline accuracy, confirming the robustness of our game-theoretic formulation.

\vspace{0.8em}

\RC{Comment 2.2: The experiments to adaptive, colluding, and stealthy Byzantine behaviors in Byzantine resilient federated learning may be added for robustness analysis.}

\AR We have added a dedicated evaluation of three additional Byzantine attack categories in the revised Section~6.3 and new Table~4a. For adaptive attacks, we implemented an adversary that observes aggregate model updates and adjusts its attack strategy every five rounds to maximize model degradation; FedGTD retains 93.1\% performance under this scenario. For colluding attacks, we configured all Byzantine clients to coordinate their gradient manipulations to approximate the honest gradient direction while injecting a persistent backdoor; FedGTD retains 92.4\% performance. For stealthy attacks, we implemented a slow-poisoning strategy where Byzantine updates remain within the honest gradient variance envelope but introduce a gradual bias; FedGTD retains 91.7\% performance. In all three scenarios, FedGTD outperforms the best baseline (RobustFL) by 4.2--6.8\%.

\vspace{0.8em}

\RC{Comment 2.3: The ablation studies may be added.}

\AR The revised manuscript now includes an expanded ablation study (Section~6.2, Table~3) that systematically evaluates the contribution of each framework component. Starting from a FedAvg baseline, we incrementally add the game-theoretic framework, domain-adaptive learning rates, Byzantine-resilient aggregation, imbalance-weighted objectives, cross-domain transfer, and adversarial training, measuring detection accuracy, Byzantine robustness, and privacy utility at each step. We also provide a new robustness-mechanism-specific ablation (Table~8) as described in our response to Reviewer~1, Comment~1.6.

\vspace{0.8em}

\RC{Comment 2.4: The paper lacks mathematical support to the proposed method. Relevant mathematical equations may be included.}

\AR We have substantially increased the mathematical rigor throughout the revised manuscript. Specifically, we have added the following new mathematical content: explicit definitions of all player strategy spaces (Definition~4.1), formal specification of the transition kernel for the stochastic game (Equation~6a in the revised text), a complete specification of the domain-adaptive loss function with all terms defined (Equation~4), the cross-domain projection mechanism expressed as a learned linear map $P_d : \mathbb{R}^{d_d} \to \mathbb{R}^{d_{\text{shared}}}$ (new Equation~4a), and the reputation update rule $r_k^{(d,t+1)} = \lambda r_k^{(d,t)} + (1-\lambda) \cos(\tilde{g}_k^{(d,t)}, \bar{g}_d^{(t)})$ with decay parameter $\lambda = 0.9$ (new Equation~13a). Every mathematical symbol is defined upon first use, and we have included a remark after each theorem explaining the practical significance of the result.

\vspace{0.8em}

\RC{Comment 2.5: The preprocessing tasks such as [four specific papers listed] may also be used to check their effect on the performance of the method.}

\AR We thank the reviewer for these suggestions. We have reviewed these works and incorporated them where relevant to our preprocessing discussion. Aggarwal~[31] is cited in the revised Section~5.1 when discussing the role of advanced feature engineering and perceptual representations in enhancing classification performance for security event data. Kumar~[32] is cited in Section~3.4 in the context of multi-sensor data fusion techniques that inform our approach to integrating heterogeneous data sources from IoT, container, and SOC environments. Aggarwal~[33] is cited in Section~5.4 when discussing the impact of data augmentation and noise on model performance, connecting to our adversarial training approach.

We note that while these works address preprocessing in domains different from network security, the underlying principles of data quality assessment and noise-robust feature engineering are applicable to our multi-domain preprocessing pipeline. We have added a brief preprocessing ablation in the revised Supplementary Material (Appendix~E.4) showing how removing individual preprocessing steps affects downstream detection accuracy across all three domains.

\vspace{1.5em}
\hrule
\vspace{1em}

%% ============================
\section*{Response to Reviewer \#3}
%% ============================

\RC{Comment 3.1: I think this paper is ready for publication. It only needs a more detailed explanation of the ideas in the formulation of Equation 1.}

\AR We thank the reviewer for this positive assessment and the specific suggestion. We have expanded the discussion surrounding Equation~1 (the main optimization objective) in Section~1.2 of the revised manuscript. The revised text now explains each component in detail.

The objective function $\sum_d \sum_{k \in \mathcal{D}_d} w_k^{(d)} \mathcal{L}_k^{(d)}(\theta_k^{(d)})$ represents the weighted aggregation of local empirical risk across all $K$ defenders and all domains. The weights $w_k^{(d)} = n_k^{(d)} / \sum_{j \in \mathcal{D}_d} n_j^{(d)}$ ensure that each client's contribution is proportional to its dataset size within its domain, preventing clients with small datasets from disproportionately influencing the global model. The local loss $\mathcal{L}_k^{(d)}$ is a domain-specific cross-entropy loss modified by the class weights $w_c^{(d)}$ to account for imbalance.

The three constraints capture the essential design requirements. The privacy constraint $\mathcal{R}_{\text{privacy}} \leq \epsilon_{\text{privacy}}$ bounds the cumulative information leakage about any individual data point, formalized through $(\epsilon, \delta)$-differential privacy. The robustness constraint $\mathcal{R}_{\text{robustness}} \leq \epsilon_{\text{robust}}$ ensures that the aggregate model deviates from the honest aggregate by at most $\epsilon_{\text{robust}}$ in $\ell_2$ norm even when a fraction $f$ of clients submit arbitrary updates. The communication constraint $\mathcal{R}_{\text{communication}} \leq B_{\text{comm}}$ bounds the total number of bits exchanged per federated round.

We also added an explanatory paragraph connecting this constrained optimization formulation to the game-theoretic setting, noting that the adversarial strategy $\pi_A$ in the robustness constraint means that the optimizer must find parameters that perform well against the worst-case adversary within the allowed corruption budget, which naturally leads to the minimax structure resolved through our Nash equilibrium analysis.

\vspace{1.5em}
\hrule
\vspace{1em}

\section*{Summary of Major Revisions}

The following list summarizes the principal changes made in the revised manuscript and supplementary material.

\begin{enumerate}
\item Expanded the game-theoretic framework with explicit game structure, player strategy spaces, utility functions, and equilibrium concept (Section~4.1, addressing Reviewers~1 and 2).
\item Added a new game-theoretic sensitivity analysis section with parametric studies and tornado diagram (Section~6.8, addressing Reviewer~2).
\item Included new experiments on adaptive, colluding, and stealthy Byzantine attacks (Section~6.3, addressing Reviewer~2).
\item Expanded ablation studies with robustness-mechanism-specific analysis and comparison against five standalone robust aggregation methods (Section~6.2, Tables~3, 8, 9, addressing Reviewers~1 and 2).
\item Provided detailed mathematical formulations for all algorithmic components with definitions of every symbol upon first use (Sections~3--4, addressing Reviewers~2 and 3).
\item Elaborated on Equation~1 with extended discussion of each component and constraint (Section~1.2, addressing Reviewer~3).
\item Described the aggregation rule, class imbalance handling, and feature heterogeneity reconciliation in detail (Section~4.2, addressing Reviewer~1).
\item Specified threat model assumptions and adversary knowledge boundaries (Section~3.2.1, addressing Reviewer~1).
\item Expanded dataset description with labeling mechanisms and cross-domain partitioning details (Section~5.1, addressing Reviewer~1).
\item Added deployment architectural details with per-component latency breakdowns (Sections~4.5 and 6.7, addressing Reviewer~1).
\item Incorporated all nine suggested references where they materially support our discussion (addressing Reviewers~1 and 2).
\item Revised writing throughout for measured academic tone, removed promotional language, and condensed redundant passages (addressing Reviewer~1).
\item Restructured Section~4 to clearly separate theoretical, algorithmic, and system engineering contributions (addressing Reviewer~1).
\item Expanded Supplementary Material with full proofs, preprocessing ablation, and extended sensitivity analysis details (Appendices~C, E.4, and new Appendix~K).
\end{enumerate}

\vspace{1em}
\noindent We believe these revisions comprehensively address all reviewer concerns. We are grateful for the opportunity to improve the manuscript and remain at your disposal for any further questions.

\vspace{1em}
\noindent Sincerely,

\vspace{0.5em}
\noindent Roger Nick Anaedevha (on behalf of all authors)

\end{document}
